\section{Limitations}\label{sec:limits}

\paragraph{Test-set reuse across versions.}
All decisions were made on validation, but the test periods were scored after every version. The test results then
shaped which ideas we tried next. For example, the CPCB losses motivated v0.6--v0.9. Later CPCB gains are therefore
exposed to adaptive overfitting to the test period. Clean evidence will come from data never scored before. To that
end we have started a prospective real-time track: live airport forecasts are logged before the weather happens and
scored afterwards.

\paragraph{Multiple comparisons.}
Each league involves dozens of tests. Verdicts with very small $p$-values, such as the WMO foundation-model wins
($p<10^{-5}$), survive a Holm correction. Borderline ones, such as the two CPCB hybrid wind wins ($p=0.02$, $0.03$), would not survive.

\paragraph{One design per league.}
The final designs differ across leagues: v0.4's design is used at the airports and the WMO stations, and v0.8's
and v0.9's on CPCB. Versions 0.6--0.9 have not yet been run as full leagues at the airports or the WMO stations.

\paragraph{Missing baseline.}
Our post-processing baseline is the bias-corrected four-model mean. A tuned gradient-boosted or neural
post-processor \citep{rasp2018postprocessing} trained on the same inputs would be a stronger opponent for the
hybrid.

\paragraph{Scale and siting.}
The leagues are small: six airports, 20 WMO stations and 25 CPCB sensors, all in north India, with a single test period
of six to nine months each. The CPCB sensors are not sited to WMO standards. Results there describe forecasts of what those sensors will
read.

\paragraph{Opponents we could not run.}
GraphCast at 0.25°, FourCastNet 3 and Aurora exceeded the memory of the free T4 GPU. GenCast and WeatherNext at full
resolution have no public archive. Our GenCast run is the small 1° model with one member, on 88 days.

\paragraph{Data dependencies.}
The archived global forecasts come from a third-party service \citep{openmeteo}. TimesFM 3.0 is used for evaluation
only, under its non-commercial licence.

\section{Conclusion}

Trained only on a network's own history, a one-million-parameter station-graph model beats TimesFM 2.5, TimesFM 3.0
and Chronos-2 on every variable at the airports and the WMO stations.
Used as a learned correction of global forecasts, the same model beats every operational global model at
north-Indian airports when all are limited to forecasts that exist at the forecast start. The hybrid does not beat AIFS's
freshest runs on airport temperature and humidity. The pure model does not beat the foundation models on wind at
Delhi's rooftop sensors, or bias-corrected FourCastNet SFNO and GenCast Mini on airport temperature; the hybrid has
not yet been scored against those two. The data threshold points to the remedy: at the airports, 24 months of local
history remove every loss to the foundation models and 48 months turn all of them into wins, so pretraining on
many more Indian stations is the natural next step. We report the full development record, including
the rejected changes, so that others can judge which gains are robust.
